
\clearpage
\section*{NeurIPS Paper Checklist}

\begin{enumerate}

\item {\bf Claims}
    \item[] Question: Do the main claims made in the abstract and introduction accurately reflect the paper's contributions and scope?
    \item[] Answer: \answerYes{}
    \item[] Justification: The abstract and \S\ref{sec:intro} state three contributions (BenchCAD dataset and CadGen pipeline; capability-decomposed task suite including the first verified parametric edit task; the Holistic Spatial and Detailing Deficit / CAD Operational Blindspot / Industrial Common Sense and Parametric Abstraction Gap diagnostics). Each claim is substantiated in \S\ref{sec:dataset}--\S\ref{sec:results} with experimental results on 10+ models.

\item {\bf Limitations}
    \item[] Question: Does the paper discuss the limitations of the work performed by the authors?
    \item[] Answer: \answerYes{}
    \item[] Justification: Appendix~\ref{app:limitations} is a dedicated Limitations section discussing the standard-parametric-vs-industrial-CAD gap, the standard-anchored-vs-compliant distinction, and the bounded scope of the Industrial Common Sense Gap measurement.

\item {\bf Theory assumptions and proofs}
    \item[] Question: For each theoretical result, does the paper provide the full set of assumptions and a complete (and correct) proof?
    \item[] Answer: \answerNA{}
    \item[] Justification: The paper presents an empirical benchmark and dataset; it makes no theoretical claims requiring proofs.

    \item {\bf Experimental result reproducibility}
    \item[] Question: Does the paper fully disclose all the information needed to reproduce the main experimental results of the paper to the extent that it affects the main claims and/or conclusions of the paper (regardless of whether the code and data are provided or not)?
    \item[] Answer: \answerYes{}
    \item[] Justification: \S\ref{sec:experiments} lists the 10+ evaluated models, decoding settings, and per-task scoring; Appendix~\ref{app:reproduce} provides the reproduction recipe (data fetch command, runner invocation, adapter wrappers, training config). Dataset, code, and baseline checkpoint are released under permissive licences.


\item {\bf Open access to data and code}
    \item[] Question: Does the paper provide open access to the data and code, with sufficient instructions to faithfully reproduce the main experimental results, as described in supplemental material?
    \item[] Answer: \answerYes{}
    \item[] Justification: All BenchCAD splits are released on Hugging Face (\texttt{BenchCAD/cad\_bench}, \texttt{BenchCAD/cad\_bench\_edit}) under CC-BY-4.0; the CadGen pipeline, evaluation harness, and BenchCAD-trained baseline checkpoint are released under MIT. Appendix~\ref{app:reproduce} gives exact commands.


\item {\bf Experimental setting/details}
    \item[] Question: Does the paper specify all the training and test details (e.g., data splits, hyperparameters, how they were chosen, type of optimizer) necessary to understand the results?
    \item[] Answer: \answerYes{}
    \item[] Justification: \S\ref{sec:exp_setup} specifies the evaluation settings; \S\ref{sec:exp_training} specifies the open BenchCAD-trained baseline's training recipe (Qwen3-VL-2B; SFT three splits + GRPO). Appendix~\ref{app:reproduce} provides the runner commands and points to the released training config.

\item {\bf Experiment statistical significance}
    \item[] Question: Does the paper report error bars suitably and correctly defined or other appropriate information about the statistical significance of the experiments?
    \item[] Answer: \answerNo{}
    \item[] Justification: All reported numbers are single-seed evaluations due to the cost of running 10+ models including proprietary APIs. We report per-family stratification (Table~\ref{tab:family_cliff}) so variance can be assessed across difficulty tiers. Future releases will include multi-seed error bars for the open BenchCAD-trained baseline.

\item {\bf Experiments compute resources}
    \item[] Question: For each experiment, does the paper provide sufficient information on the computer resources (type of compute workers, memory, time of execution) needed to reproduce the experiments?
    \item[] Answer: \answerYes{}
    \item[] Justification: \S\ref{sec:exp_setup} describes inference compute; \S\ref{sec:exp_training} describes baseline training; Appendix~\ref{app:reproduce} provides the runner invocation.

\item {\bf Code of ethics}
    \item[] Question: Does the research conducted in the paper conform, in every respect, with the NeurIPS Code of Ethics?
    \item[] Answer: \answerYes{}
    \item[] Justification: The dataset is procedurally generated, contains no PII or proprietary designs, and references public engineering standards by code only. No human-subjects data is collected.


\item {\bf Broader impacts}
    \item[] Question: Does the paper discuss both potential positive societal impacts and negative societal impacts of the work performed?
    \item[] Answer: \answerYes{}
    \item[] Justification: Appendix~\ref{app:ethics} discusses the dataset's lack of personal or proprietary content, the potential dual-use risk inherent to any improved CAD-capable model, and the absence of negative impacts specific to BenchCAD beyond those of the broader open-source CAD ecosystem.

\item {\bf Safeguards}
    \item[] Question: Does the paper describe safeguards that have been put in place for responsible release of data or models that have a high risk for misuse?
    \item[] Answer: \answerNA{}
    \item[] Justification: BenchCAD contains expert-authored parametric CAD parts referencing public engineering standards. The released open baseline (Qwen3-VL-2B) poses no high-risk misuse potential beyond what is already available in open-source CAD ecosystems.

\item {\bf Licenses for existing assets}
    \item[] Question: Are the creators or original owners of assets used in the paper properly credited and are the license and terms of use explicitly mentioned and properly respected?
    \item[] Answer: \answerYes{}
    \item[] Justification: We credit DeepCAD~\citep{Wu2021}, Fusion360 Gallery~\citep{Willis2021fusion}, Qwen2-VL~\citep{Qwen2VL}, and the cadrille/CADEvolve checkpoints with explicit licence statements. Re-distribution of derived data complies with the source licences.

\item {\bf New assets}
    \item[] Question: Are new assets introduced in the paper well documented and is the documentation provided alongside the assets?
    \item[] Answer: \answerYes{}
    \item[] Justification: BenchCAD ships with a Croissant 1.0 metadata file (validated by the public Croissant checker), per-family parameter schemas, dataset card, and the datasheet of Appendix~\ref{app:datasheet}. The CadGen pipeline and evaluation harness include README, environment specification, and reproduction commands.

\item {\bf Crowdsourcing and research with human subjects}
    \item[] Question: For crowdsourcing experiments and research with human subjects, does the paper include the full text of instructions given to participants and screenshots, if applicable, as well as details about compensation (if any)?
    \item[] Answer: \answerYes{}
    \item[] Justification: BenchCAD does not involve human-subjects evaluation; the dataset is procedurally generated and all reported metrics are computed by automated execution and verification pipelines.

\item {\bf Institutional review board (IRB) approvals or equivalent for research with human subjects}
    \item[] Question: Does the paper describe potential risks incurred by study participants, whether such risks were disclosed to the subjects, and whether Institutional Review Board (IRB) approvals (or an equivalent approval/review based on the requirements of your country or institution) were obtained?
    \item[] Answer: \answerNA{}
    \item[] Justification: The only human-subjects component is paid professional CAD engineers scoring CAD samples on a standard consulting basis; no personal data is collected and no risk to participants beyond ordinary professional work. IRB review is not required by the authors' institutions for this engagement type.

\item {\bf Declaration of LLM usage}
    \item[] Question: Does the paper describe the usage of LLMs if it is an important, original, or non-standard component of the core methods in this research?
    \item[] Answer: \answerYes{}
    \item[] Justification: LLMs are the evaluated subjects, not core method components. Where LLMs assist in dataset construction (e.g., GPT-5-mini in CADEvolve's evolutionary expansion, used here only as a transfer baseline), this is described in \S\ref{sec:exp_training}. The CadGen procedural pipeline contains no LLM in the loop. LLM use in writing assistance is incidental.

\end{enumerate}
